% ============================================================
%  Appendix.tex — Phụ lục
% ============================================================
\section*{PHỤ LỤC}
\phantomsection
\addcontentsline{toc}{section}{\numberline{}PHỤ LỤC}
\setcounter{section}{0}
\renewcommand{\thesection}{\Alph{section}}

% ============================================================
\section{Mã nguồn IDS (Contiki-NG)}
\label{app:code}
% ============================================================

\subsection{Module phát hiện DIS Flooding}

% TODO: dán mã nguồn C vào đây, ví dụ:
% \lstinputlisting[language=C,
%   caption={ids\_dis\_flooding.c}]{Code/ids_dis_flooding.c}

\begin{lstlisting}[language=C, caption={Skeleton: ids\_dis\_flooding.c}]
#include "contiki.h"
#include "net/routing/rpl-lite/rpl.h"

/* Nguong so ban tin DIS trong mot chu ky Trickle */
#define DIS_FLOOD_THRESHOLD  10

static uint16_t dis_count = 0;

void ids_dis_on_receive(void) {
  dis_count++;
  if (dis_count > DIS_FLOOD_THRESHOLD) {
    /* TODO: Kich hoat canh bao DIS Flooding */
  }
}
\end{lstlisting}

\subsection{Module phát hiện Blackhole}

\begin{lstlisting}[language=C, caption={Skeleton: ids\_blackhole.c}]
#include "contiki.h"

#define PDR_THRESHOLD  0.7   /* PDR duoi 70% -> nghi ngo Blackhole */

float ids_compute_pdr(uint16_t sent, uint16_t received) {
  if (sent == 0) return 1.0f;
  return (float)received / (float)sent;
}
\end{lstlisting}

\subsection{Module phát hiện Decreased Rank}

\begin{lstlisting}[language=C, caption={Skeleton: ids\_decreased\_rank.c}]
#include "contiki.h"
#include "net/routing/rpl-lite/rpl.h"

void ids_rank_check(rpl_dag_t *dag, rpl_parent_t *parent,
                    rpl_rank_t advertised_rank) {
  rpl_rank_t expected = /* tinh toan rank du kien */ 0;
  if (advertised_rank < expected) {
    /* TODO: Danh dau parent nay la nghi ngo */
  }
}
\end{lstlisting}

\subsection{Module phát hiện VeRA / Version Number Attack}

\begin{lstlisting}[language=C, caption={Skeleton: ids\_vera.c}]
#include "contiki.h"
#include "net/routing/rpl-lite/rpl.h"

static uint8_t last_known_version = 0;

void ids_version_check(uint8_t received_version) {
  if (received_version > last_known_version + 1) {
    /* TODO: Version tang bat thuong -> canh bao VeRA attack */
  }
  last_known_version = received_version;
}
\end{lstlisting}

% ============================================================
\section{Cấu hình mô phỏng Cooja}
\label{app:cooja}
% ============================================================

\subsection{Tham số mô phỏng}

\begin{table}[H]
\centering
\caption{Tham số cấu hình chung trong Cooja}
\begin{tabularx}{\textwidth}{lX}
  \toprule
  \textbf{Tham số} & \textbf{Giá trị} \\
  \midrule
  Mô phỏng radio    & Unit Disk Graph Medium (UDGM) \\
  Số nút tổng       & 20 nút (1 Border Router + 19 mote) \\
  Tỉ lệ nút tấn công & 0\%, 10\%, 25\%, 50\% \\
  Thời gian mô phỏng & 300 giây / kịch bản \\
  Tần suất gửi dữ liệu & 1 gói / 10 giây \\
  Hệ điều hành      & Contiki-NG 4.x \\
  Objective Function & MRHOF (ETX metric) \\
  \bottomrule
\end{tabularx}
\end{table}

\subsection{Hướng dẫn chạy mô phỏng}

\begin{enumerate}
  \item Cài đặt Contiki-NG và Cooja theo hướng dẫn tại
        \url{https://github.com/contiki-ng/contiki-ng/wiki}.
  \item Mở file \texttt{simulation/rpl\_ids\_scenario.csc} trong Cooja.
  \item Chọn kịch bản tấn công trong menu \textbf{Simulation Parameters}.
  \item Nhấn \textbf{Start} và quan sát log từ \textbf{Mote Output}.
  \item Xuất kết quả ra file \texttt{.csv} để phân tích với script Python
        trong thư mục \texttt{analysis/}.
\end{enumerate}
